\chapter{Boolean Analysis — Block Characteristic}
%%% generated by the blueprint pipeline from TCSlib.BooleanAnalysis.GaussianSpace.Noise.BlockCharacteristic
\section{Overview}

Declarations of the Lean module \texttt{TCSlib.BooleanAnalysis.GaussianSpace.Noise.BlockCharacteristic}.

\section{Declarations}

\begin{definition}[Correlated sign characteristic]
\lean{BooleanAnalysis.GaussianSpace.correlatedSignCharacteristic}
\label{BooleanAnalysis.GaussianSpace.correlatedSignCharacteristic}
\leanok
For $
ho,a,b\in\mathbb R$, define the correlated sign characteristic by
$\frac{1+\rho}{2}\cos(a+b)+\frac{1-\rho}{2}\cos(a-b)$.
\end{definition}

\begin{lemma}[Cosine identity using the sinc function]
\lean{BooleanAnalysis.GaussianSpace.cos_sub_one_eq_sinc}
\label{BooleanAnalysis.GaussianSpace.cos_sub_one_eq_sinc}
\leanok
\difficulty{4}
For every real number $t$, $\cos t-1=-\frac{t^2}{2}\,\operatorname{sinc}(t/2)^2$.
\end{lemma}

\begin{lemma}[Cosine limit under square-root scaling]
\lean{BooleanAnalysis.GaussianSpace.tendsto_nat_mul_cos_div_sqrt_sub_one}
\label{BooleanAnalysis.GaussianSpace.tendsto_nat_mul_cos_div_sqrt_sub_one}
\leanok
\uses{BooleanAnalysis.GaussianSpace.cos_sub_one_eq_sinc}
\difficulty{5}
For every real number $a$, the sequence $m\bigl(\cos(a/\sqrt m)-1\bigr)$, indexed by
natural numbers $m$, converges as $m\to\infty$ to $-a^2/2$.
\end{lemma}

\begin{lemma}[Second-order limit of the correlated sign characteristic]
\lean{BooleanAnalysis.GaussianSpace.tendsto_nat_mul_correlatedSignCharacteristic_sub_one}
\label{BooleanAnalysis.GaussianSpace.tendsto_nat_mul_correlatedSignCharacteristic_sub_one}
\leanok
\uses{BooleanAnalysis.GaussianSpace.correlatedSignCharacteristic, BooleanAnalysis.GaussianSpace.tendsto_nat_mul_cos_div_sqrt_sub_one}
\difficulty{3}
For all real numbers $\rho,a,b$, as $m\to\infty$ through positive integers,
$m\bigl(C_\rho(a/\sqrt m,b/\sqrt m)-1\bigr)\to-\tfrac12(a^2+b^2+2\rho ab)$, where
$C_\rho(s,t)=\frac{1+\rho}{2}\cos(s+t)+\frac{1-\rho}{2}\cos(s-t)$.
\end{lemma}

\begin{lemma}[Gaussian limit of correlated sign characteristics]
\lean{BooleanAnalysis.GaussianSpace.tendsto_correlatedSignCharacteristic_pow}
\label{BooleanAnalysis.GaussianSpace.tendsto_correlatedSignCharacteristic_pow}
\leanok
\uses{BooleanAnalysis.GaussianSpace.correlatedSignCharacteristic, BooleanAnalysis.GaussianSpace.tendsto_nat_mul_correlatedSignCharacteristic_sub_one}
\difficulty{3}
For real numbers $\rho,a,b$, the sequence
\[
\left(\frac{1+\rho}{2}\cos\left(\frac{a+b}{\sqrt
m}\right)+\frac{1-\rho}{2}\cos\left(\frac{a-b}{\sqrt m}\right)\right)^m
\]
converges as the natural number $m$ tends to infinity to
$\exp\!\left(-\frac{a^2+b^2+2\rho ab}{2}\right)$.
\end{lemma}

\begin{lemma}[Characteristic function of a correlated Gaussian pair]
\lean{BooleanAnalysis.GaussianSpace.integral_exp_correlatedGaussianMeasure}
\label{BooleanAnalysis.GaussianSpace.integral_exp_correlatedGaussianMeasure}
\leanok
\uses{BooleanAnalysis.GaussianSpace.correlatedGaussianMeasure, BooleanAnalysis.GaussianSpace.integral_exp_mul_I_gaussianPi}
\difficulty{6}
Let $n$ be a nonnegative integer, let $\rho\in[-1,1]$, and let $a,b\in\mathbb R^n$. For
a pair $(X,Y)$ with the correlated Gaussian law obtained from independent standard
Gaussian vectors by setting $Y_i=\rho X_i+\sqrt{1-\rho^2}\,Z_i$, the characteristic
function of $\sum_i(a_iX_i+b_iY_i)$ is $\mathbb
E[e^{i\sum_i(a_iX_i+b_iY_i)}]=\exp\!\left(-\frac12\sum_i(a_i^2+b_i^2+2\rho a_i
b_i)\right)$.
\end{lemma}
